\section{强证据实际带来了什么}

2024 年的 Communications Biology 文献之所以重要，是因为它支持这样一个主张：
意识状态伴随着有组织的动力学标志。
对于本手稿而言，这主要是一种治理结果：本书可以在不显得生物学上幼稚的前提下，讨论有结构的状态组织。

2025 年的 Communications Biology 文献又增加了另一层价值。
通过把皮层连接性、局部动力学与全局意识状态的稳定性相关项连接起来，
它支持了这样一种观点：状态持续既依赖耦合架构，也依赖局部系统行为。
从业务意义上说，这座桥梁把日常的“当前心情”式表达，升级为更严肃的“稳定性架构”词汇。

Nature Neuroscience 关于睡眠转变的文献，是最有用的转变锚点，
因为它处理的是一个具体的生物学机制切换，而不是关于努力的道德叙事。
这使得分岔式语言与本书后文关于决策窗口和杠杆点的讨论真正相关。